% Chapter 1

\chapter{绪论}

\section{研究背景与意义}

随着人工智能生成内容技术的快速发展，基于扩散模型\cite{ref1}的图像生成方法已经趋于成熟，研究重心也逐步从静态图像扩展到视频生成。相较于图像生成，视频生成不仅要求模型具有较强的画面细节表达能力，具备单帧画面的细节质量和艺术表现力，还要求生成结果在时间维度上保持连续、稳定和自然，因此其建模难度显著高于图像生成，这使得视频生成成为当前生成式人工智能领域的重要研究方向\cite{ref2}。尤其是在动画制作、短视频创作、数字文旅展示、游戏资产生产和虚拟内容传播等场景中，若能够以较低门槛生成具有明确艺术风格的视频内容，将具有显著的应用价值。

目前主流视频生成模型虽然在通用场景下已取得较好效果，但在面对特定风格定制化任务时仍存在两个突出问题。其一，视频扩散模型参数规模较大，若直接采用全参数微调，不仅训练成本高，而且对显存、训练时长和数据规模要求都较为苛刻，普通研究环境和消费级硬件难以长期支撑；其二，风格迁移往往容易破坏视频原有的时序一致性，使生成结果出现闪烁、动作变形或主体不稳定等现象，导致生成视频难以兼顾画面风格和连续观感。因此，如何在有限算力条件下实现高质量、可控且时序连贯的特定风格视频生成，成为一个值得深入研究的问题。

LoRA\cite{ref3}作为一种典型的参数高效微调方法，通过冻结预训练底模的大部分参数，仅在关键模块如线性层中引入少量可训练的低秩矩阵，从而显著降低训练成本与部署门槛。在图像生成领域，LoRA已被广泛用于风格迁移、主体定制和轻量微调等任务\cite{ref4}；若将这一思路延伸到视频生成场景，一方面可以减少模型适配过程中的资源消耗，另一方面也为探索空间风格建模与时间动态建模之间的平衡提供了新的技术路径。围绕这一思路开展研究，有助于形成一种低成本、高效率的特定风格视频生成方案。

结合本课题的实际条件与研究目标，本文不以从零训练大型视频基础模型为目标，而是围绕“低成本完成特定风格学习并验证其向视频阶段迁移的可行性”展开研究。通过构建特定风格数据集、设计基于扩散模型的LoRA微调方案、比较不同底模先验的影响，并将图像阶段获得的风格能力扩展到短视频原型验证，本文力图回答一个具有现实意义的问题：在有限算力和有限数据条件下，是否能够建立一条兼顾风格表达、实验可复现性与工程可实施性的特定风格视频生成路线。

从理论层面看，本课题有助于分析LoRA在视频扩散模型中的作用机制，特别是其对空间注意力层和时间注意力层的不同影响，从而为参数高效微调方法在高维时空生成任务中的应用提供参考。从实践层面看，本课题致力于在较低算力条件下实现特定风格的视频生成流程，有望为独立动画创作、个性化内容生产和中小团队的智能视频制作提供可复用的技术方案，具有较强的现实意义。

\section{国内外研究现状}

国外对视频生成的研究起步较早，技术路线经历了从生成对抗网络向扩散模型的明显转移。早期方法多关注短时序、低分辨率视频的生成\cite{ref5}，随着扩散模型在图像任务中表现出更稳定的训练特性和更高的生成质量\cite{ref6}，研究重心逐步转向基于扩散过程的视频生成框架。近年来，Stable Video Diffusion\cite{ref7}、Lumiere\cite{ref8}等代表性工作在长时序建模、画面清晰度提升和时空联合结构设计等方面取得了较大进展，为后续风格迁移和参数高效适配研究提供了更加稳固的基础。

在参数高效微调方面，LoRA最初广泛应用于大语言模型领域\cite{ref3}，随后被快速迁移到视觉生成任务中\cite{ref9,ref10}。国外研究者围绕Few-Shot和One-Shot场景展开了大量探索，尝试在较少样本条件下完成风格迁移、主体保持和个性化生成等任务。相关研究表明，LoRA在静态图像风格微调中具有显著优势，但在视频任务中，如何处理时间注意力层、如何在保持动作连贯性的同时注入特定风格，仍然是尚未完全解决的问题\cite{ref11}。

国内研究更强调视频生成技术的工程落地与可控生成能力的提升。以AnimateDiff\cite{ref12}、VideoCrafter\cite{ref13}等为代表的工作，一方面从运动模块建模、视频高分辨率生成与时序先验复用等方向推动了视频生成技术的发展，另一方面也逐渐形成了风格模块与运动模块解耦的工程化思路。一些研究通过将风格模块与运动模块解耦，增强了视频生成的灵活性；另一些研究则引入ControlNet\cite{ref14}、ReferenceNet\cite{ref15}等控制机制，用于提升角色一致性、结构可控性和动作稳定性。这些成果表明，国内在视频生成系统化实现、开源复现与应用适配方面已经积累了较为扎实的经验。

综合国内外研究现状可以看出，当前特定风格视频生成大致呈现出三方面趋势。第一，模型适配方式正由粗粒度统一微调转向更细粒度的模块化控制，即分别考虑空间外观建模与时间动态建模\cite{ref16}。第二，风格生成不再局限于单一适配器，而逐步发展为多LoRA组合、结构约束与参考图引导等联合方案，以同时兼顾画质、风格强度和主体稳定性\cite{ref17,ref18}。第三，研究重点从能否生成视频转向如何生成可控、可复现且更符合目标风格的视频，这意味着评价体系和实验流程也在不断完善\cite{ref19,ref20}。

尽管已有研究为本课题提供了较多基础，但仍存在两个与本文高度相关的空白。其一，面向小规模特定风格数据集的轻量化实验流程仍缺乏统一范式\cite{ref21}，尤其是在学生研究环境中，往往很难兼顾训练成本与结果可解释性。其二，许多工作更关注通用视频生成能力，而对图像域风格学习如何向视频阶段迁移讨论不足\cite{ref22}。本文正是在这一背景下展开，希望通过图像阶段的LoRA风格学习、底模替换比较和短视频原型验证，为特定风格视频生成提供一条更具可操作性的研究路线。

\section{论文研究内容}

围绕基于LoRA技术的特定风格视频生成模型这一主题，本文首先从理论与方法层面对扩散模型\cite{ref1}、潜空间生成\cite{ref4}、参数高效微调\cite{ref3}以及视频时序建模\cite{ref12,ref13}相关技术进行梳理，在此基础上分析LoRA适用于视频生成任务的原因，并明确其在训练成本、风格迁移能力和可部署性方面的优势。通过理论分析与相关技术比较，为后续实验方案设计提供支撑。

在数据层面，本文围绕目标风格构建了较为规范的数据处理流程。图像阶段对风格样本进行整理、筛选、裁剪、尺寸统一和文本描述对齐\cite{ref23}，形成可直接用于LoRA训练的图文配对数据；视频阶段则对片段样本进行帧序列导出、样本划分和一致性检查，为后续短视频原型实验提供基础。与单纯堆叠素材不同，本文更强调训练数据在风格表达、语义描述\cite{ref24}和实验可复现性之间的平衡。

在模型与实验层面，本文以Stable Diffusion系列模型为基础构建LoRA风格学习方案，重点研究LoRA在U-Net注意力模块\cite{ref25,ref26}中的注入方式，以及训练步数、学习率、LoRA秩、caption组织方式等因素对生成结果的影响。除以Stable Diffusion1.5为主线完成短训与长训验证外，本文还进一步引入StorybookRedmond和Counterfeit-V3.0等具有不同先验风格的底模进行替换实验\cite{ref27}，以分析底模先验对目标风格迁移强度和内容保真度的影响。

在结果验证层面，本文采用先图像、后视频的渐进式研究路线。图像阶段通过单图推理、批量提示词对比、差异可视化和统计指标分析\cite{ref28}来判断LoRA是否成功学习到目标风格；视频阶段则在图像LoRA的基础上构建递推式短视频生成原型，并从相邻帧差异、边缘稳定性和亮度波动\cite{ref29,ref30}等角度评估风格能力向连续帧迁移的效果。该路线既符合本课题的现实算力条件，也能够形成较为完整的实验闭环。

归纳而言，本文的研究内容可以概括为四个方面：一是完成特定风格视频生成相关理论与技术的梳理；二是构建面向风格学习的数据处理与训练流程；三是设计并验证基于LoRA的多底模风格微调方案；四是通过图像和短视频实验分析其有效性、局限性及进一步优化方向。上述内容共同构成了全文的研究主线。

\section{论文章节安排}

全文共分为五章。第一章为绪论，主要介绍研究背景与意义、国内外研究现状、论文研究内容及整体结构安排。第二章为相关理论与关键技术基础，重点阐述扩散模型、潜空间扩散框架、LoRA参数高效微调、视觉—语言标注技术、视频生成与结构控制方法等内容。第三章为系统设计与实现，围绕数据集构建、预处理流程、LoRA训练方案、图像推理与短视频原型生成等模块展开说明。第四章为实验与结果分析，对多轮图像训练、底模替换实验和视频原型验证结果进行定量与定性分析。第五章为总结与展望，对全文工作进行归纳，并讨论当前研究的不足及后续可继续深入的方向。

\section{本章小结}

本章围绕论文研究工作的总体背景进行了展开，分析了特定风格视频生成任务在当前生成式人工智能发展背景下的研究价值，梳理了国内外相关研究的主要进展与发展趋势，并明确了本文的研究内容与章节安排。通过本章的论述，可以看出，基于LoRA的参数高效微调方法为解决特定风格视频生成中的高成本训练与时序一致性问题提供了可行思路，也为后续相关理论分析、系统设计与实验验证奠定了整体框架。
